% Dirichlet Round 10 German TeX
% Work: G. Lejeune Dirichlet, Werke II, pp. 139--158.
% Checked against Werke II scan images 152--171.
\documentclass[11pt]{article}
\usepackage[a4paper,margin=1in]{geometry}
\usepackage[T1]{fontenc}
\usepackage[utf8]{inputenc}
\usepackage{lmodern}
\usepackage{amsmath,amssymb}
\usepackage{microtype}
\setlength{\parindent}{1.2em}
\setlength{\parskip}{0.25em}
\newcommand{\cf}[1]{(#1)}
\begin{document}

\begin{center}
{\Large G. Lejeune Dirichlet's Werke. II.}\par\vspace{1.2em}
{\Large XIII.}\par\vspace{0.5em}
{\LARGE \textbf{Vereinfachung der Theorie\ der binären quadratischen Formen\ von positiver Determinante.}}\par\vspace{1em}
{\large Von G. Lejeune Dirichlet.}\par\vspace{0.5em}
{\small Abhandlungen der Königl. Preussischen Akademie der Wissenschaften von 1854, S. 99--115.}
\end{center}

\begin{center}
{\large \textbf{Vereinfachung der Theorie der binären quadratischen Formen\ von positiver Determinante.}}\par
{\small Gelesen in der Akademie der Wissenschaften am 13. Juli 1854.\footnote{Die Mittheilung im Akademie-Bericht von 1854, S. 384 lautet: ``Hr. Dirichlet las: Ueber eine Eigenschaft der Kettenbrüche und deren Gebrauch zur Vereinfachung der Theorie der quadratischen Formen von positiver Determinante.'' K.}}
\end{center}

Je grösser der Umfang ist, welchen die höhere Arithmetik durch das Epoche machende Werk von Gauss und andere spätere Arbeiten gewonnen hat, um so wünschenswerther erscheint es, dass der Zugang zu diesem schönen Zweige der Analysis durch Vereinfachung des elementaren Theiles desselben so viel als möglich erleichtert werde. In solcher Absicht habe ich schon in mehreren früheren Abhandlungen meinen Untersuchungen die dazu erforderlichen bekannten Sätze mit neuer Begründung vorausgeschickt: eine ähnliche Vereinfachung bezweckt der gegenwärtige Aufsatz, welcher der Theorie der quadratischen Formen von positiver Determinante gewidmet ist. Bekanntlich erfordert diese Lehre in ihrer bisherigen Gestalt sehr ins Einzelne gehende Betrachtungen, die sich, wie die folgende Darstellung zeigen wird, daraus entfernen lassen.

Ich beginne mit einigen Bemerkungen über Kettenbrüche, die, obgleich ihrem wesentlichen Inhalte nach nicht neu, in der für die hier davon zu machende Anwendung geeigneten Form vorauszuschicken sind.

\section*{\S. 1.}

Ein endlicher oder unendlicher Kettenbruch wie
\[
  \alpha+\frac{1}{\beta+\frac{1}{\gamma+\cdots}}
\]
soll im Folgenden durch
\[
  (\alpha,\beta,\gamma,\ldots)
\]
bezeichnet werden. Wir haben nur Kettenbrüche zu betrachten, deren sämmtliche Glieder ganze Zahlen sind, natürlich mit Ausnahme des letzten für den Fall, wo die Entwickelung nicht zu Ende geführt ist, und wo dieses Glied als ein sogenannter vollständiger Quotient jeden anderen Werth haben kann. Von besonderer Wichtigkeit für arithmetische Untersuchungen sind diejenigen Kettenbrüche, deren Glieder bis auf das erste, für welches auch der Werth Null zulässig ist, positiv sind. Durch einen solchen Kettenbruch lässt sich eine positive Irrationalgrösse $\omega$ nur auf eine Weise ausdrücken; und wir wollen die Darstellung von $\omega$ in dieser Form, oder, wenn $\omega$ negativ ist, die Darstellung ihres absoluten Werthes mit vorgesetztem negativen Zeichen, die normale Kettenbruch-Entwickelung von $\omega$ nennen.

Wir haben nun zunächst die Aufgabe zu behandeln, aus einem Kettenbruche wie
\[
  \omega=(\alpha,\beta,\ldots,\mu,\nu,p,q,r,\ldots,u,v,\ldots),
\]
in welchem die Glieder erst von $p$ ab, und zwar $p$ eingeschlossen, sämmtlich positiv sind, die normale Entwickelung der Irrationalgrösse $\omega$ abzuleiten. Es wird sich leicht zeigen lassen, dass dies durch eine Reihe von Umformungen bewerkstelligt werden kann, bei welchen die Glieder, die auf ein hinlänglich entferntes $u$ folgen, unberührt bleiben, und dass die Anzahl der neuen Glieder, welche schliesslich an die Stelle von $\alpha,\beta,\ldots,\nu$ getreten sind, von der Anzahl der letzteren um eine gerade oder ungerade Zahl verschieden sein wird, je nachdem $\omega$ positiv oder negativ ist.

Man betrachte zunächst den Fall, wo $\nu$ nicht das erste Glied ist. Unter dieser Voraussetzung kann man $\mu,\nu$ und einige der unmittelbar folgenden Glieder, während alle übrigen ungeändert bleiben, durch neue Glieder ersetzen, deren Anzahl von der Anzahl jener um eine gerade Zahl verschieden ist, und welche mit Ausnahme des ersten, welches Null oder negativ sein kann, sämmtlich positiv sind. Bei dieser partiellen Umformung hat man zu unterscheiden, ob $\nu$ Null ist oder einen negativen Werth $-n$ hat. Im ersteren Falle sind die drei Glieder $\mu,0,p$ durch das einzige Glied $\mu+p$ zu ersetzen. Der andere Fall zerfällt in die drei Unterabtheilungen
\[
 n>1,\qquad n=1,\ p>1,\qquad n=1,\ p=1,
\]
denen entsprechend eine der folgenden Gleichungen anzuwenden ist:
\[
 (\mu,-n,p,q,\ldots)=(\mu-1,1,n-2,1,p-1,q,\ldots),
\]
\[
 (\mu,-1,p,q,\ldots)=(\mu-2,1,p-2,q,\ldots),
\]
\[
 (\mu,-1,1,q,r,s,\ldots)=(\mu-q-2,1,r-1,s,\ldots).
\]
Die durch eine solche partielle Umformung hervorgebrachte Änderung der Gliederzahl beträgt beziehungsweise $2,0,-2$ Einheiten. Wenn eine der Differenzen
\[
 n-2,\quad p-1,
\quad p-2,
\quad r-1
\]
sich auf Null reducirt, so ist für die Null und die beiden benachbarten positiven Glieder ein einziges, der Summe der letzteren gleiches Glied zu setzen.\footnote{Dass sich die negativen Glieder aus einem Kettenbruche entfernen lassen, hat schon Lagrange bemerkt; aber die von ihm zu diesem Zwecke gegebene Gleichung, welche mit der ersten der obigen zusammenfällt, reicht nicht aus, da sie für den Fall $p=1$ ein neues negatives Glied einführt. Will man dieses durch abermalige Anwendung derselben Gleichung beseitigen, so wird man zu dem ursprünglichen Kettenbruche zurückgeführt. Oeuvres de Lagrange, t. 2 p. 624.}

Durch wiederholte Anwendung desselben Verfahrens lässt sich bewirken, dass alle Glieder, vom zweiten einschliesslich ab, positiv werden. Ist dann zugleich das erste nicht negativ, so ist die Operation geschlossen. Hat hingegen das erste Glied einen negativen Werth $-a$ und folglich der Kettenbruch die Form
\[
  \omega=(-a,b,c,d,\ldots),
\]
so hat man, je nachdem $b>1$ oder $b=1$ ist,
\[
 \omega=-(a-1,1,b-1,c,\ldots),
 \qquad\hbox{oder}\qquad
 \omega=-(a-1,c+1,d,\ldots)
\]
zu setzen. Das Resultat stimmt wieder mit dem früher Behaupteten überein.

\section*{\S. 2.}

I. Finden zwischen zwei Grössen $\omega,\Omega$ und den ganzen Zahlen $\alpha,\beta,\gamma,\delta$, deren erste nicht Null ist, die Relationen
\[
 \omega=\frac{\gamma+\delta\Omega}{\alpha+\beta\Omega},
 \qquad
 \alpha\delta-\beta\gamma=1,
\]
statt, so lässt sich immer eine Gleichung der Form
\[
  \omega=(\lambda,m,\ldots,r,\sigma,\Omega)
\]
bilden, in welcher von den ganzen Zahlen $\lambda,m,\ldots,r,\sigma$ nur die erste und letzte Null oder negativ sein können, die Zwischenglieder aber, wenn sie nicht ganz fehlen, positiv und in gerader Anzahl sind.

Da man nach der Form der vorausgesetzten Gleichungen die Zeichen von $\alpha,\beta,\gamma,\delta$ gleichzeitig ändern kann, so darf $\alpha$ positiv angenommen werden. Ist nun $\alpha=1$, so hat man sogleich
\[
 \omega=\frac{\gamma+(\beta\gamma+1)\Omega}{1+\beta\Omega}
       =(\gamma,\beta,\Omega).
\]
Ist hingegen $\alpha>1$, so verwandle man $\gamma/\alpha$ auf die gewöhnliche Weise in einen Kettenbruch, indem man alle Divisionsreste positiv wählt. Man erhält
\[
 \frac{\gamma}{\alpha}=(\lambda,m,\ldots,r),
\]
in welchem nur $\lambda$ Null oder negativ sein kann, und die Anzahl der Glieder $m,\ldots,r$ gerade vorausgesetzt werden kann, da sich das Glied $r$, für welches man zunächst einen Werth, der grösser als Eins ist, erhält, nöthigenfalls in $(r-1,1)$ auflösen lässt. Da die zu diesem Kettenbruche gehörigen Näherungsbrüche
\[
 \frac{\lambda}{1},\quad \frac{\lambda m+1}{m},\quad\ldots,\quad
 \frac{\varphi}{f},\quad \frac{\gamma}{\alpha}
\]
irreductibel sind und positive Nenner haben, so wird der letzte derselben auch der Form nach mit $\gamma/\alpha$ zusammenfallen. Nach dem bekannten Satze
\[
 \alpha\varphi-\gamma f=1
\]
ergiebt der Vergleich mit $\alpha\delta-\beta\gamma=1$ die Darstellung
\[
  \beta=\alpha\sigma+f,
  \qquad
  \delta=\gamma\sigma+\varphi,
\]
wo $\sigma$ eine ganze Zahl ist. Der Bruch $\delta/\beta$ lässt sich also vermittelst des neuen Gliedes $\sigma$ der Reihe der Näherungsbrüche anschliessen, und man hat
\[
  \omega=(\lambda,m,\ldots,r,\sigma,
\Omega).
\]

II. Für das Folgende ist noch der besondere Fall näher zu betrachten, wo $\alpha,\beta,\gamma,\delta$ sämmtlich positiv sind und zugleich die Bedingungen
\[
 \gamma\geq \alpha,
 \qquad
 \delta>\gamma
\]
erfüllen. Dann sind $\lambda$ und $\sigma$ positiv. Ist $\alpha=1$, so liegt dies schon in der Voraussetzung, denn in diesem Falle ist $\lambda=\gamma$, $\sigma=\beta$. Ist dagegen $\alpha>1$, so ist $\lambda$, welches der unmittelbar unter $\gamma/\alpha$ liegenden ganzen Zahl gleich ist, positiv. Auch $\sigma$ ist positiv; denn aus
\[
 \delta=\gamma\sigma+\varphi
\]
würde bei $\sigma=0$ die Ungleichung $\delta=\varphi<\gamma$, und bei negativem $\sigma$ sogar ein negativer Werth von $\delta$ folgen. Bezeichnen wir die positiven Zahlen $\lambda,\sigma$ der Gleichförmigkeit wegen mit $l,s$, so ist in diesem besonderen Falle
\[
 \frac{\gamma}{\alpha}=(l,m,\ldots,r),
 \qquad
 \frac{\delta}{\beta}=(l,m,\ldots,r,s),
 \qquad
 \omega=(l,m,\ldots,r,s,\Omega),
\]
wo die Glieder $l,m,\ldots,r,s$ sämmtlich positiv und in gerader Anzahl sind.

\section*{\S. 3.}

Wir wenden uns jetzt zu den quadratischen Formen
\[
  ax^2+2bxy+cy^2=(a,b,c),
\]
welche alle dieselbe positive Determinante
\[
  D=b^2-ac
\]
haben sollen. Die positive ganze Zahl $D$ ist beliebig, nur darf sie kein Quadrat sein. Da hiernach die äusseren Coefficienten $a,c$ nie Null sind, so ist eine Form durch $D$, den mittleren Coefficienten und einen der äusseren Coefficienten vollständig bestimmt.

Jeder Form $(a,b,c)$ lassen wir die Gleichung
\[
 a+2b\omega+c\omega^2=0
\]
entsprechen, deren Wurzeln
\[
 \omega=\frac{-b\mp\sqrt D}{c}
\]
immer so dargestellt werden sollen, dass der unveränderte dritte Coefficient $c$ den Nenner bildet. Dem oberen und unteren Zeichen entsprechend werden die beiden Werthe als erste und zweite zur Form gehörige Wurzeln unterschieden. Eine Form ist durch ihre Determinante und eine der zu ihr gehörigen Wurzeln völlig bestimmt. Gehört nämlich derselbe Werth zu beiden Formen $(a,b,c)$, $(A,B,C)$ als erste Wurzel oder zu beiden als zweite, so hat man
\[
 \frac{-b\mp\sqrt D}{c}=\frac{-B\mp\sqrt D}{C},
\]
und wegen der Irrationalität von $\sqrt D$ folgt $B=b$, $C=c$, also die Identität der beiden Formen.

Wenn zwei Formen
\[
 ax^2+2bxy+cy^2,
 \qquad
 AX^2+2BXY+CY^2
\]
äquivalent genannt werden, so ist stets die eigentliche Äquivalenz gemeint: es existiert eine Substitution
\[
 x=\alpha X+\beta Y,
 \qquad
 y=\gamma X+\delta Y,
\]
deren Coefficienten die Bedingung
\[
 \alpha\delta-\beta\gamma=1
\]
erfüllen und durch welche die erste Form in die zweite übergeht. Substitutionen mit Determinante $-1$ werden im Folgenden ausgeschlossen.

I. Zwischen den gleichnamigen zu äquivalenten Formen gehörigen Wurzeln $\omega$ und $\Omega$ und den Coefficienten der Substitution besteht immer die Gleichung
\[
 \omega=\frac{\gamma+\delta\Omega}{\alpha+\beta\Omega}.
\]
Bringt man diese Gleichung in die Form
\[
 \Omega=\frac{\gamma-\alpha\omega}{\beta\omega-\delta},
\]
setzt für $\omega$ seinen Werth ein und befreit den Nenner von der Irrationalität, so wird die rechte Seite
\[
 \frac{-M\mp\sqrt D}{N},
\]
wo
\[
 M=a\alpha\beta+b(\alpha\delta+\beta\gamma)+c\gamma\delta,
\qquad
 N=a\beta^2+2b\beta\delta+c\delta^2.
\]
Diese Ausdrücke $M,N$ sind gerade die Coefficienten $B,C$ der transformierten Form; also ist die Behauptung erwiesen.

II. Findet umgekehrt die eben angegebene Gleichung für ein Paar gleichnamiger Wurzeln statt und erfüllen die ganzen Zahlen $\alpha,\beta,\gamma,\delta$ die Bedingung $\alpha\delta-\beta\gamma=1$, so sind die Formen äquivalent und die erste geht durch die angegebene Substitution in die zweite über. Denn ohne neue Rechnung erhält man
\[
 \frac{-B\mp\sqrt D}{C}=\frac{-M\mp\sqrt D}{N},
\]
folglich $B=M$, $C=N$.

III. Häufig werden benachbarte Formen auftreten, d. h. Formen
\[
 (a,b,a'),
 \qquad
 (a',b',a'')
\]
bei denen der letzte Coefficient der ersten mit dem ersten der zweiten zusammenfällt und deren mittlere Coefficienten zugleich die Bedingung
\[
 b+b'\equiv0\pmod {a'}
\]
erfüllen. Solche Formen sind immer äquivalent. Wendet man auf die erste die Substitution
\[
 \begin{pmatrix}0&1\\-1&\delta\end{pmatrix}
\]
an, so erhält man eine neue Form mit erstem Coefficienten $a'$ und mit zweitem Coefficienten $-b-a'\delta$, der dem gegebenen $b'$ gleich wird, wenn
\[
 \delta=-\frac{b+b'}{a'}
\]
gesetzt wird. Für die zugehörigen gleichnamigen Wurzeln $\omega,\omega'$ gilt
\[
 \omega=\delta-\frac1{\omega'},
 \qquad
 \omega'=-\frac1{\omega-\delta}.
\]

\section*{\S. 4.}

Wenn von den beiden zur Form $(a,b,c)$ gehörigen Wurzeln
\[
 \frac{-b-\sqrt D}{c},
 \qquad
 \frac{-b+\sqrt D}{c}
\]
die erste ihrem absoluten Werthe nach über, die zweite unter der Einheit liegt, und diese Wurzeln überdies entgegengesetzte Zeichen haben, so heisst die Form reducirt. Aus der ersten Bedingung folgt $b>0$, aus der zweiten $b<\sqrt D$. Das Product $-ac=D-b^2$ ist positiv; die äusseren Coefficienten $a,c$ haben also entgegengesetzte Zeichen, und das Zeichen der ersten Wurzel stimmt mit dem von $a$ überein und ist dem Zeichen von $c$ entgegengesetzt.

Ist $(a,b,c)$ reducirt, so ist auch $(c,b,a)$ reducirt. Denn jede Wurzel der einen Form ist der reciproke Werth der ungleichnamigen Wurzel der anderen.

Für jede Determinante $D$ giebt es nur eine endliche Anzahl von reducirten Formen. Man erhält sie, indem man für jedes positive $b<\sqrt D$ alle positiven und negativen Factoren $c$ von $D-b^2$ aufsucht, deren absolute Werthe zwischen $\sqrt D+b$ und $\sqrt D-b$ liegen, und dann
\[
 a=-\frac{D-b^2}{c}
\]
setzt.

Sei nun $(a,b,a')$ eine gegebene reducirte Form. Unter den rechts benachbarten Formen $(a',b',a'')$, deren mittlere Coefficienten durch
\[
 b+b'\equiv0\pmod {a'}
\]
bestimmt werden, giebt es genau eine reducirte. Soll $(a',b',a'')$ reducirt sein, so muss die erste Wurzel $\omega'$ dieser Form dem Zeichen nach der ersten Wurzel $\omega$ von $(a,b,a')$ entgegengesetzt sein. In
\[
 \omega'=-\frac1{\omega-\delta}
\]
ist also die ganze Zahl $\delta$ so zu wählen, dass $\omega'$ ein unechter Bruch werde und im Zeichen $\omega$ entgegengesetzt sei. Gleichbedeutend ist die Forderung, dass $\omega-\delta$ ein echter Bruch werde und im Zeichen mit $\omega$ übereinstimme. Dies geschieht auf eine und nur eine Art: man wählt für $\delta$ diejenige der beiden $\omega$ unmittelbar benachbarten ganzen Zahlen, welche auf derselben Seite von $\omega$ liegt wie die Null. Da $|\omega|>1$, ist dieses $\delta$ nicht Null, hat dasselbe Zeichen wie $\omega$, und liegt seinem absoluten Werthe nach unmittelbar unter $|\omega|$.

Dass die so bestimmte Nachbarform wirklich reducirt ist, sieht man, wenn in derselben Gleichung $\omega$ die zweite Wurzel bedeutet. Dann hat $\omega'$ das Vorzeichen von $-\delta$; zugleich ist $\omega-\delta$ ein unechter Bruch, also $\omega'$ ein echter Bruch mit dem Zeichen von $\delta$, also entgegengesetzt zum Zeichen der ersten Wurzel der neuen Form.

Setzt man, für die erste Wurzel $\omega$,
\[
 \omega-\delta=\sigma,
\]
wobei $\sigma$ ein echter Bruch vom selben Zeichen wie $\omega$ ist, so folgt aus
\[
 b'=-b-a'\delta
\]
nach Einsetzung des Werthes von $\omega$:
\[
 b'=\sqrt D+a'\sigma.
\]
Da $\sigma$ das Zeichen von $\omega$ hat und folglich das entgegengesetzte von $a'$, liegt $b'$ zwischen $\sqrt D$ und $\sqrt D\mp a'$, wo das obere oder untere Zeichen gilt, je nachdem $a'$ positiv oder negativ ist. Zusammen mit der Congruenz
\[
 b'\equiv -b\pmod {a'}
\]
bestimmt dies $b'$ leicht und ohne Zweideutigkeit. Ebenso giebt es eine und nur eine links benachbarte reducirte Form $({}'a,{}'b,a)$.

\section*{\S. 5.}

Bildet man aus einer reducirten Form $\varphi_0$ die ihr rechts benachbarte $\varphi_1$, aus dieser wieder $\varphi_2$ und so fort, und verfährt ebenso nach links, so erhält man eine nach beiden Seiten unendliche Reihe äquivalenter reducirter Formen
\[
 \ldots,\varphi_{-2},\varphi_{-1},\varphi_0,
 \varphi_1,\varphi_2,\ldots.
\]
Da es nur endlich viele reducirte Formen von gegebener Determinante giebt, wiederholen sich die in dieser Reihe enthaltenen Formen. Da die ersten Coefficienten abwechselnd positiv und negativ sind, können zwei Formen nur dann identisch sein, wenn die Differenz ihrer Indices gerade ist. Ferner folgt aus der Bildungsweise, nach welcher jedes Glied sein vorhergehendes und folgendes vollständig bestimmt, dass, wenn zwei Formen identisch sind, auch alle gleich weit nach derselben Seite von ihnen abstehenden Formen identisch sind.

Sei unter den auf $\varphi_0$ folgenden Formen $\varphi_{2n}$ die erste mit $\varphi_0$ identische. Dann sind
\[
 \varphi_0,\varphi_1,\varphi_2,\ldots,\varphi_{2n-1}
\]
alle verschieden. Diese $2n$ Formen bilden eine Periode, die sich nach beiden Seiten wiederholt; zwei Formen $\varphi_\mu,\varphi_\nu$ sind genau dann identisch, wenn
\[
 \mu\equiv \nu\pmod {2n}.
\]
Da eine Form und die zu ihr gehörigen gleichnamigen Wurzeln sich gegenseitig bestimmen, gilt dieselbe Bedingung für die Gleichheit gleichnamiger Wurzeln $\omega_\mu,\omega_\nu$.

Bezeichne $\delta_\nu$ die im absoluten Werthe der ersten Wurzel $\omega_\nu$ enthaltene grösste ganze Zahl, mit dem Zeichen von $\omega_\nu$ genommen. Zwischen gleichnamigen Wurzeln gilt
\[
 \omega_\nu=\delta_\nu-
\frac1{\omega_{\nu+1}}.
\]
Die Congruenz $\mu\equiv\nu\pmod{2n}$ zieht $\delta_\mu=\delta_\nu$ nach sich, aber nicht umgekehrt.

Wir wählen den Index so, dass die ersten Coefficienten der Formen mit geradem Index positiv sind. Dann stimmen die Zeichen von $\omega_\nu$ und $\delta_\nu$ mit $(-1)^\nu$ überein. Setzt man
\[
 \delta_\nu=(-1)^\nu k_\nu,
\]
so sind $k_\nu$ positive ganze Zahlen, für die, wegen $\delta_{\nu+2n}=\delta_\nu$,
\[
 k_{\nu+2n}=k_\nu
\]
gilt. Aus den Gleichungen für die Wurzeln folgt durch alternirende Multiplication mit $+1$ und $-1$ die normale periodische Entwickelung
\[
 \omega_\nu=(-1)^\nu(k_\nu,k_{\nu+1},\ldots,k_{2n+
u-1},k_\nu,k_{\nu+1},\ldots).
\]
Ferner erhält man für die zweite Wurzel $\omega'_
u$
\[
 \frac1{\omega'_
u}=(-1)^{\nu-1}(k_{\nu-1},k_{\nu-2},\ldots,k_{\nu-2n},k_{\nu-1},k_{\nu-2},\ldots).
\]
Die Periode für den reciproken Werth der zweiten Wurzel besteht daher aus den Gliedern der ersten Periode in umgekehrter Ordnung.

Die Periode $k_0,k_1,
\ldots,k_{2n-1}$ hat keine kürzere gerade Periode. Gäbe es eine solche mit $2m$ Gliedern, so würden aus der Entwickelung der ersten Wurzel $\omega_0$ die Wurzeln $\omega_0$ und $\omega_{2m}$, und folglich die Formen $\varphi_0$ und $\varphi_{2m}$, zusammenfallen, entgegen der Wahl von $2n$.

Endlich kann die Gesammtheit der reducirten Formen einer gegebenen Determinante immer in solche Perioden vertheilt werden. Man beginnt mit einer reducirten Form und bildet die Periode, der sie angehört; bleiben reducirte Formen übrig, so verfährt man mit einer von ihnen ebenso, bis alle erschöpft sind.

\section*{\S. 6.}

Wir kommen nun zur Entscheidung, ob zwei gegebene Formen äquivalent sind. Da sich jede Form leicht auf eine äquivalente reducirte Form bringen lässt, und Formen derselben Periode äquivalent sind, bleibt nur zu untersuchen, ob Formen aus verschiedenen Perioden äquivalent sein können.

Wir wählen in den beiden zu vergleichenden Perioden Formen mit positivem erstem Coefficienten und schreiben
\[
 \varphi_0=(a,b,c),
 \qquad
 \Phi_0=(A,B,C),
\]
wobei die zugehörigen ersten Wurzeln durch die normalen Kettenbrüche
\[
 \omega_0=(k_0,k_1,k_2,\ldots),
 \qquad
 \Omega_0=(K_0,K_1,K_2,\ldots)
\]
dargestellt seien. Nehmen wir an, die Formen seien äquivalent und die erste gehe durch eine Substitution mit Determinante $1$ in die zweite über. Dann gilt nach \S. 3
\[
 \omega_0=\frac{\gamma+\delta\Omega_0}{\alpha+\beta\Omega_0}.
\]
Hier kann $\alpha$ nicht Null sein; sonst wäre $\gamma=\pm1$ und folglich $A=c$, was der Vorzeichenvoraussetzung widerspricht. Nach \S. 2 besitzt man also eine Darstellung
\[
 \omega_0=(\lambda,m,\ldots,r,\sigma,
 K_0,K_1,K_2,\ldots).
\]
Die Zahl der eingeschobenen Glieder $\lambda,m,\ldots,r,\sigma$ sei $2g$. Wird dieser Kettenbruch nach \S. 1 in die normale Form gebracht, so ändert sich, weil $\omega_0$ positiv ist, die Anzahl der Glieder bis zu einem hinlänglich entfernten unverändert bleibenden Glied $K_\nu$ um eine gerade Zahl $2h$. Für alle genügend grossen $\nu$ muss daher
\[
 K_\nu=k_{2g+2h+\nu}
\]
gelten. Schreibt man statt $\nu$ die Zahl $2N i+
u$, wobei $2Ni$ ein hinlänglich grosses Vielfaches der Periode $2N$ der $K$-Reihe ist, und reducirt den Index $2g+2h+2Ni$ der $k$-Reihe modulo $2n$ auf $2m$, so erhält man für alle $\nu\geq0$
\[
 K_\nu=k_{2m+
u}.
\]
Folglich ist
\[
 \Omega_0=\omega_{2m},
 \qquad
 \Phi_0=\varphi_{2m}.
\]
Die zweite Form liegt also in der zur ersten gehörigen Periode. Formen aus verschiedenen Perioden können demnach nicht äquivalent sein.

\section*{\S. 7.}

Nachdem der schwierigste Satz der Theorie auf diese Weise bewiesen ist, bleibt nur anzudeuten, wie die übrige Lehre in den Punkten zu modificiren ist, die nicht allein auf diesem Satze, sondern auf seiner Begründung beruhen.

Da die Operationen, durch welche man die Äquivalenz zweier Formen erkennt, immer eine erste Substitution ergeben, durch welche die eine Form in die andere übergeht, bleibt noch die Aufgabe, aus einer gegebenen Transformation einer Form in eine andere alle übrigen abzuleiten. Diese Aufgabe führt man leicht auf die Darstellung aller Transformationen einer Form in sich selbst zurück. Dabei darf vorausgesetzt werden, dass die Coefficienten keinen gemeinschaftlichen Theiler haben; denn jede Substitution, durch welche eine Form in sich selbst übergeht, wirkt auch bei der durch Entfernung des gemeinschaftlichen Theilers entstehenden Form und umgekehrt.

Ist also $(a,b,c)$ eine primitive Form und ist $\sigma$ der positive grösste gemeinschaftliche Theiler von $a,2b,c$, so ist $\sigma=1$ oder $2$; der letztere Fall kann nur eintreten, wenn $D=b^2-ac$ von der Form $4h+1$ ist. Dann werden alle Substitutionen, welche die Form in sich selbst verwandeln, durch
\[
 \alpha=\frac{t-bu}{\sigma},
 \qquad
 \beta=-\frac{cu}{\sigma},
 \qquad
 \gamma=\frac{au}{\sigma},
 \qquad
 \delta=\frac{t+bu}{\sigma}
\]
erhalten, wenn $t,u$ alle ganzen Lösungen der Gleichung
\[
 t^2-Du^2=\sigma^2
\]
durchlaufen.\footnote{Vgl. Disquisitiones arithmeticae, p. 181, oder Crelle's Journal, Band 24, S. 328; der dort für complexe Zahlen gegebene Beweis bleibt wörtlich für reelle Zahlen anwendbar, wenn das dortige Zeichen $\omega$ hier als $\sigma$ verstanden wird.}

Für eine reducirte Form lässt sich das Transformationsproblem direct lösen. Wir dürfen $a>0$ voraussetzen und uns auf diejenigen Substitutionen beschränken, deren Coefficienten $\alpha,\beta,\gamma,\delta$ sämmtlich positiv sind. Sei
\[
 \omega=(k_0,k_1,\ldots,k_{2n-1},k_0,k_1,\ldots)
\]
die normale periodische Entwickelung der ersten Wurzel. Bezeichnen $\gamma/\alpha$ und $\delta/\beta$ zwei auf einander folgende Näherungswerthe, deren zweiter dem Endgliede $k_{2n-1}$ irgend einer Periode entspricht, so hat man
\[
 \alpha\delta-\beta\gamma=1,
 \qquad
 \omega=\frac{\gamma+\delta\omega}{\alpha+\beta\omega}.
\]
Nach \S. 3 geht die Form durch die aus vier positiven Coefficienten gebildete Substitution in sich selbst über.

Umgekehrt erhält man alle positiven Substitutionen so. Sind $\alpha,\beta,\gamma,\delta$ die Coefficienten einer solchen Substitution, so gilt
\[
 \omega=\frac{\gamma+\delta\omega}{\alpha+\beta\omega},
\]
also
\[
 \beta\omega^2+(\alpha-\delta)\omega-
\gamma=0.
\]
Da die erste Wurzel zwischen $1$ und $\infty$, die zweite zwischen $-1$ und $0$ liegt, ist die linke Seite für $\omega=1$ negativ und für $\omega=-1$ positiv. Daraus folgen die Ungleichheiten
\[
 \gamma-\alpha>\beta-
\delta,
 \qquad
 \delta-\gamma>\alpha-
\beta,
\]
und weiter
\[
 \gamma\geq\alpha,
 \qquad
 \delta>\gamma.
\]
Damit sind alle Voraussetzungen von \S. 2, II erfüllt, und es gilt
\[
 \frac{\gamma}{\alpha}=(l,m,\ldots,r),
 \qquad
 \frac{\delta}{\beta}=(l,m,\ldots,r,s),
 \qquad
 \omega=(l,m,\ldots,r,s,\omega).
\]
Setzt man die periodische Entwickelung von $\omega$ ein, so sind die beiden normalen Kettenbrüche identisch; folglich besteht die Reihe $l,m,\ldots,r,s$ aus einer oder mehreren Perioden $k_0,k_1,\ldots,k_{2n-1}$, und $\gamma/\alpha$, $\delta/\beta$ sind auf einander folgende Näherungswerthe, deren zweiter dem Ende einer Periode entspricht. Die kleinsten positiven Substitutionscoefficienten entsprechen dem Ende der ersten Periode; sie werden aus den kleinsten positiven Werthen von $t$ und $u$ erhalten.

In der That folgt aus den Formeln für $\alpha$ und $\delta$
\[
 \alpha\delta=\frac{t^2-b^2u^2}{\sigma^2}
 =\frac{t^2-Du^2}{\sigma^2}-\frac{ac u^2}{\sigma^2}
 =1-
\frac{ac u^2}{\sigma^2}.
\]
Da $-ac>0$, haben $\alpha$ und $\delta$ dasselbe Zeichen, welches wegen
\[
 \alpha+\delta=\frac{2t}{\sigma}
\]
mit dem Zeichen von $t$ übereinstimmt. Ebenso folgt aus den Ausdrücken für $\beta$ und $\gamma$, dass auch diese, wenn sie nicht beide Null sind, was $u=0$ voraussetzt, im Zeichen mit $u$ übereinstimmen. Die oben untersuchten positiven Substitutionen werden also aus positiven $t,u$ erhalten; da $\beta$ und $\gamma$ mit $u$ wachsen, entspricht die in den kleinsten Zahlen ausgedrückte Substitution den kleinsten positiven Werthen von $t,u$, die, sobald jene aus dem Kettenbruche bestimmt ist, durch
\[
 t=\frac{\sigma}{2}(\delta+\alpha),
 \qquad
 u=\frac{\sigma}{2b}(\delta-
\alpha)
\]
gefunden werden.

\end{document}
